% =====================================================================
%  Rapport — Plateforme d'aide à la décision multicritère (MCDM)
%  Module « Aide à la décision » — Prof. Lamrani Alaoui, EMI Rabat
%  Compilation : pdflatex rapport.tex  (deux passes pour la table des matières)
% =====================================================================
\documentclass[11pt,a4paper]{report}

% ---------- Encodage et langue ----------
\usepackage[utf8]{inputenc}
\usepackage[T1]{fontenc}
\usepackage[french]{babel}
\usepackage{lmodern}
\usepackage{microtype}

% ---------- Mise en page ----------
\usepackage[a4paper,margin=2.5cm,headheight=14pt]{geometry}
\usepackage{fancyhdr}
\usepackage{titlesec}
\usepackage{setspace}
\usepackage{enumitem}
\setlist{itemsep=2pt,topsep=4pt}
\onehalfspacing

% ---------- Mathématiques ----------
\usepackage{amsmath,amssymb,amsthm}
\usepackage{bm}

% ---------- Tableaux, figures, couleurs ----------
\usepackage{booktabs}
\usepackage{array}
\usepackage{tabularx}
\usepackage{multirow}
\usepackage{float}
\usepackage{caption}
\usepackage[table]{xcolor}
\usepackage{graphicx}
\usepackage{tikz}
\usetikzlibrary{arrows.meta,positioning,shapes.geometric,fit,babel}

% ---------- Code ----------
\usepackage{listings}

% ---------- Liens ----------
\usepackage[hidelinks]{hyperref}

% ---------- Paramètres à compléter ----------
\newcommand{\auteurs}{\textit{[Nom Prénom de l'étudiant(e)]}}
\newcommand{\filiere}{\textit{[Filière / option]}}
\newcommand{\annee}{2025--2026}

% ---------- Couleurs et styles ----------
\definecolor{emiblue}{HTML}{1C5CAB}
\definecolor{lightblue}{HTML}{CDE2FB}
\definecolor{codebg}{HTML}{F6F6F4}
\definecolor{codecomment}{HTML}{52514E}
\definecolor{okgreen}{HTML}{006300}
\definecolor{warnorange}{HTML}{B35A00}

\newcommand{\ok}{\textcolor{okgreen}{\textbf{conforme}}}
\newcommand{\ecart}{\textcolor{warnorange}{\textbf{écart}}}

\lstdefinestyle{py}{
  language=Python,
  basicstyle=\ttfamily\small,
  backgroundcolor=\color{codebg},
  keywordstyle=\color{emiblue}\bfseries,
  commentstyle=\color{codecomment}\itshape,
  stringstyle=\color{warnorange},
  showstringspaces=false,
  frame=single, rulecolor=\color{lightblue},
  breaklines=true, columns=fullflexible,
  numbers=left, numberstyle=\tiny\color{codecomment},
  xleftmargin=1.5em, framexleftmargin=1.5em
}
\lstset{style=py}

\titleformat{\chapter}[display]
  {\normalfont\huge\bfseries\color{emiblue}}{\chaptertitlename\ \thechapter}{12pt}{\Huge}
\titlespacing*{\chapter}{0pt}{-20pt}{24pt}

\pagestyle{fancy}
\fancyhf{}
\fancyhead[L]{\small\textcolor{codecomment}{Plateforme MCDM}}
\fancyhead[R]{\small\textcolor{codecomment}{\leftmark}}
\fancyfoot[C]{\thepage}
\renewcommand{\headrulewidth}{0.4pt}

\theoremstyle{definition}
\newtheorem{definition}{Définition}[chapter]
\theoremstyle{remark}
\newtheorem*{remarque}{Remarque}

\renewcommand{\arraystretch}{1.2}
\renewcommand{\lstlistingname}{Code}
\renewcommand{\lstlistlistingname}{Liste des codes}

% =====================================================================
\begin{document}

% ---------------------------------------------------------------------
% Page de garde
% ---------------------------------------------------------------------
\begin{titlepage}
  \centering
  {\large Université Mohammed V de Rabat\par}
  {\Large\bfseries École Mohammadia d'Ingénieurs (EMI)\par}
  \vspace{0.6cm}
  {\large Module : Aide à la décision\par}
  \vspace{2.5cm}
  {\color{emiblue}\rule{\linewidth}{1.2pt}\par}
  \vspace{0.4cm}
  {\Huge\bfseries Conception et réalisation d'une plateforme d'aide à la décision multicritère (MCDM)\par}
  \vspace{0.4cm}
  {\Large Pondération, classement, évaluation absolue, MCDM floue et analyse de robustesse sous Streamlit\par}
  \vspace{0.4cm}
  {\color{emiblue}\rule{\linewidth}{1.2pt}\par}
  \vspace{2.5cm}
  \begin{tabular}{rl}
    \textbf{Réalisé par :} & \auteurs \\
    \textbf{Filière :} & \filiere \\
    \textbf{Encadré par :} & Prof. Lamrani Alaoui \\
  \end{tabular}
  \vfill
  {\large Année universitaire \annee\par}
\end{titlepage}

% ---------------------------------------------------------------------
\pagenumbering{roman}
\chapter*{Résumé}
\addcontentsline{toc}{chapter}{Résumé}

Ce rapport présente la conception, l'implémentation et la validation d'une plateforme web d'aide à la
décision multicritère (MCDM) développée en Python avec Streamlit. La plateforme reprend fidèlement les
formulations mathématiques et la taxonomie du cours \og Aide à la décision \fg{} de l'EMI : définition de
l'espace de décision (alternatives, critères positifs et négatifs), pondération des critères par des
méthodes subjectives (AHP, BWM), objectives (Entropie, CRITIC), hybrides ou manuelles, classement des
alternatives par des méthodes compensatoires (WSM, WPM, WASPAS, TOPSIS, VIKOR) et de surclassement
(PROMETHEE II), évaluation absolue d'une alternative isolée (MAUT, seuil de décision, courbe ROC, indice de
Youden, distance à l'idéal), MCDM floue (TOPSIS flou de Chen et TOPSIS flou intuitionniste avec
l'opérateur IFWA) et analyse de robustesse (perturbation des poids, simulation de Monte Carlo, corrélation
de rang de Spearman, consensus de Borda).

Dix exercices du cours sont intégrés comme études de cas interactives et résolus pas à pas. Les résultats
calculés par la plateforme sont confrontés aux valeurs publiées dans les diapositives : ils sont conformes
pour l'ensemble des cas, à une exception près (BWM, exemple du fournisseur), pour laquelle nous montrons que
les poids affichés dans le support ne satisfont pas les contraintes du modèle linéaire présenté, alors que
la valeur optimale $\xi^*\approx 0{,}08$ est bien retrouvée. Une suite de quinze tests automatisés garantit
la reproductibilité de ces résultats.

\medskip
\noindent\textbf{Mots-clés :} aide à la décision multicritère, MCDM, AHP, BWM, CRITIC, entropie, TOPSIS,
VIKOR, WASPAS, PROMETHEE, MAUT, courbe ROC, logique floue, TOPSIS flou, ensembles flous intuitionnistes,
analyse de sensibilité, corrélation de Spearman, Streamlit.

\tableofcontents
\listoftables
\listoffigures

\clearpage
\pagenumbering{arabic}

% =====================================================================
\chapter{Introduction générale}
% =====================================================================

\section{Contexte}

Dans de nombreuses situations, un décideur doit choisir la meilleure option parmi plusieurs alternatives :
implanter un hôpital dans l'une de quatre villes candidates, électrifier un village isolé par une solution
solaire, éolienne ou hybride, retenir un fournisseur parmi cinq offres ou sélectionner un projet
d'investissement. Avec un seul critère, il suffirait de retenir l'alternative la mieux notée. En pratique,
le choix dépend de plusieurs critères \emph{contradictoires} --- rentabilité, risque, coût, impact social
ou environnemental --- et aucune alternative n'est généralement la meilleure sur tous les critères. Il faut
alors pondérer les critères et rechercher une solution de compromis : c'est le rôle des méthodes d'analyse
multicritère (\emph{Multi-Criteria Decision Making}, MCDM).

\section{Problématique}

Le cours présente une large palette de méthodes, chacune reposant sur des hypothèses différentes
(compensation ou non entre critères, nature exacte ou imprécise des données, nombre de décideurs, nombre
d'alternatives). Les appliquer à la main est long et sujet aux erreurs de calcul ; comparer plusieurs
méthodes et étudier la stabilité d'un classement face à une variation des poids l'est encore davantage.
D'où la nécessité d'un outil qui :
\begin{itemize}
  \item implémente rigoureusement les formulations du cours ;
  \item expose chaque étape intermédiaire pour garder une valeur pédagogique ;
  \item permet de comparer les méthodes et d'évaluer la robustesse des décisions.
\end{itemize}

\section{Objectifs}

La plateforme réalisée doit permettre à l'utilisateur de :
\begin{enumerate}
  \item définir l'espace de décision : alternatives, critères et sens de préférence (critère positif
        \emph{benefit} [+] ou négatif \emph{cost} [$-$]) ;
  \item choisir une méthode de pondération des critères : subjective, objective, hybride ou manuelle ;
  \item choisir la portée de l'évaluation : une seule alternative (évaluation absolue) ou plusieurs
        alternatives (classement ou surclassement) ;
  \item exécuter des méthodes MCDM classiques (\emph{crisp}) ou floues ;
  \item mener une analyse de sensibilité et une comparaison des classements entre méthodes ;
  \item rejouer de manière interactive les études de cas du cours.
\end{enumerate}

\section{Organisation du rapport}

Le chapitre~\ref{ch:theorie} rappelle le cadre théorique du MCDM. Les chapitres~\ref{ch:ponderation}
à~\ref{ch:robustesse} détaillent les méthodes implémentées, avec leurs formules exactes et nos choix
d'implémentation. Le chapitre~\ref{ch:architecture} décrit l'architecture logicielle et le
chapitre~\ref{ch:interface} l'interface. Le chapitre~\ref{ch:validation} confronte les résultats de la
plateforme aux exercices du cours et le chapitre~\ref{ch:comparaison} illustre l'analyse comparative et de
sensibilité sur un cas complet. Nous terminons par les limites, les perspectives et une conclusion.

% =====================================================================
\chapter{Cadre théorique de l'analyse multicritère}\label{ch:theorie}
% =====================================================================

\section{MODM et MADM}

On distingue deux grandes familles de méthodes MCDM :
\begin{itemize}
  \item l'\textbf{optimisation multi-objectifs} (MODM), qui optimise simultanément
        $F(x)=(f_1(x),\dots,f_k(x))$ sous des contraintes $g_i(x)\le 0$ sur un espace de décision continu ;
  \item l'\textbf{analyse multi-attributs} (MADM), qui classe et sélectionne la meilleure alternative parmi
        un ensemble \emph{fini} de $m$ alternatives $A_1,\dots,A_m$ évaluées selon $n$ critères
        $C_1,\dots,C_n$.
\end{itemize}
La plateforme traite le cas MADM. Le problème est formulé par une matrice de décision et un vecteur de
poids :
\begin{equation}
  D=\begin{pmatrix}
    x_{11} & x_{12} & \cdots & x_{1n}\\
    x_{21} & x_{22} & \cdots & x_{2n}\\
    \vdots & \vdots & \ddots & \vdots\\
    x_{m1} & x_{m2} & \cdots & x_{mn}
  \end{pmatrix},\qquad W=(w_1,\dots,w_n),\quad \sum_{j=1}^n w_j=1,
\end{equation}
où $x_{ij}$ est le score de l'alternative $A_i$ selon le critère $C_j$.

\section{Typologie des critères}

\begin{itemize}
  \item \textbf{Sens de préférence} : critères positifs (\emph{benefit}, à maximiser : profit,
        productivité) et négatifs (\emph{cost}, à minimiser : coût, délai, perte). Nous notons $J$ et $J'$
        ces deux ensembles.
  \item \textbf{Relation entre critères} : indépendants ou dépendants (ANP nécessaire dans ce dernier cas).
  \item \textbf{Nature de la mesure} : quantitatifs, ou qualitatifs exprimés par des variables
        linguistiques à convertir en échelle numérique.
\end{itemize}

\section{Taxonomie retenue}\label{sec:taxonomie}

La plateforme est organisée selon les quatre axes de la taxonomie du cours (tableau~\ref{tab:taxonomie}).

\begin{table}[H]
\centering
\caption{Taxonomie MADM et couverture par la plateforme}\label{tab:taxonomie}
\begin{tabularx}{\linewidth}{>{\raggedright\arraybackslash}p{3.4cm}X>{\raggedright\arraybackslash}p{4.2cm}}
\toprule
\textbf{Axe} & \textbf{Modalités} & \textbf{Page de la plateforme}\\
\midrule
Précision des évaluations & \emph{Crisp} (valeurs exactes) / \emph{Fuzzy} (jugements imprécis) & 3 · Évaluation / 4 · MCDM floue\\
Nombre de décideurs & \emph{Single} / \emph{Group} (agrégation des avis) & 4 · MCDM floue (K décideurs)\\
Nombre d'alternatives & Une seule (seuil, référence) / plusieurs (classement) & 3 · Évaluation (choix de la portée)\\
Étape du processus & Pondération (subjective / objective) puis classement (additif / surclassement) & 2 · Pondération / 3 · Évaluation\\
\bottomrule
\end{tabularx}
\end{table}

\section{Choix d'une méthode et bonnes pratiques}

Aucune méthode n'est meilleure que les autres dans l'absolu. Le choix dépend :
\begin{itemize}
  \item de la \textbf{nature des données} : valeurs exactes $\rightarrow$ méthodes \emph{crisp} ;
        jugements imprécis ou linguistiques $\rightarrow$ méthodes floues ;
  \item du \textbf{nombre de décideurs} : un expert $\rightarrow$ \emph{single} ; plusieurs
        $\rightarrow$ \emph{group} MCDM avec agrégation ;
  \item de la \textbf{compensation entre critères} : acceptée $\rightarrow$ méthodes additives (WSM,
        TOPSIS) ; refusée (sécurité, environnement) $\rightarrow$ surclassement (ELECTRE, PROMETHEE) ;
  \item de l'\textbf{effort demandé au décideur} : AHP exige $n(n-1)/2$ comparaisons, BWM $2n-3$, CRITIC
        et l'entropie aucune.
\end{itemize}
Le cours recommande en outre d'appliquer \emph{plusieurs} méthodes, de comparer les classements, de mener
une analyse de sensibilité et de justifier explicitement la méthode retenue. La plateforme a été conçue
pour rendre ces bonnes pratiques immédiates (chapitre~\ref{ch:robustesse}).

\begin{figure}[H]
\centering
\begin{tikzpicture}[
  node distance=0.55cm and 0.5cm,
  box/.style={draw=emiblue, fill=lightblue!50, rounded corners=3pt, align=center,
              minimum height=1.25cm, text width=2.55cm, font=\small},
  arr/.style={-{Stealth[length=2.5mm]}, thick, emiblue}]
  \node[box] (p1) {1. Identifier\\le problème};
  \node[box, right=of p1] (p2) {2. Élaborer les\\alternatives};
  \node[box, right=of p2] (p3) {3. Critères et\\collecte des données};
  \node[box, right=of p3] (p4) {4. Pondérer et\\classer};
  \node[box, below=of p4] (p5) {5. Sensibilité\\et robustesse};
  \node[box, left=of p5] (p6) {6. Décision\\justifiée};
  \draw[arr] (p1) -- (p2);
  \draw[arr] (p2) -- (p3);
  \draw[arr] (p3) -- (p4);
  \draw[arr] (p4) -- (p5);
  \draw[arr] (p5) -- (p6);
\end{tikzpicture}
\caption{Processus de prise de décision suivi par la plateforme}\label{fig:processus}
\end{figure}

% =====================================================================
\chapter{Pondération des critères}\label{ch:ponderation}
% =====================================================================

Deux familles de méthodes sont distinguées selon l'origine de l'information :
\begin{itemize}
  \item les \textbf{méthodes subjectives} (AHP, BWM) tirent les poids du jugement des experts ; elles
        reflètent les priorités réelles du décideur mais dépendent de sa cohérence ;
  \item les \textbf{méthodes objectives} (entropie, CRITIC) calculent les poids à partir de la matrice de
        décision : plus un critère discrimine les alternatives, plus son poids est élevé ; elles sont
        reproductibles mais ignorent l'importance stratégique.
\end{itemize}
En pratique, on les combine souvent en \textbf{poids hybrides}.

\section{AHP (Analytic Hierarchy Process)}

\subsection{Matrice de comparaison par paires}
Chaque critère est comparé aux autres sur l'échelle de Saaty de 1 (importance égale) à 9 (extrêmement plus
important). La matrice $A=(\alpha_{ij})$ est réciproque :
\begin{equation}
  \alpha_{ii}=1,\qquad \alpha_{ji}=\frac{1}{\alpha_{ij}}.
\end{equation}
Les poids sont le vecteur propre normalisé associé à la plus grande valeur propre : $AW=\lambda_{max}W$.

\subsection{Méthode approximative du cours}
\begin{enumerate}
  \item Calculer la somme de chaque colonne : $s_j=\sum_i \alpha_{ij}$.
  \item Diviser chaque élément par la somme de sa colonne : $\bar\alpha_{ij}=\alpha_{ij}/s_j$.
  \item Le poids $w_i$ est la moyenne de la ligne $i$ : $w_i=\frac{1}{n}\sum_j \bar\alpha_{ij}$.
\end{enumerate}

\subsection{Test de consistance}
\begin{equation}
  \lambda_{max}=\sum_{i=1}^n s_i\,w_i,\qquad CI=\frac{\lambda_{max}-n}{n-1},\qquad CR=\frac{CI}{RI},
\end{equation}
avec l'indice aléatoire $RI$ du tableau~\ref{tab:ri}. La matrice est jugée consistante si $CR<0{,}1$.

\begin{table}[H]
\centering
\caption{Indice de consistance aléatoire RI}\label{tab:ri}
\begin{tabular}{l*{8}{c}}
\toprule
$n$ & 3 & 4 & 5 & 6 & 7 & 8 & 9 & 10\\
\midrule
RI & 0,58 & 0,90 & 1,12 & 1,24 & 1,32 & 1,41 & 1,45 & 1,56\\
\bottomrule
\end{tabular}
\end{table}

\subsection{AHP pour le classement}
Pour chaque critère, une matrice de comparaison des alternatives fournit des \emph{priorités locales}
$w_{ji}$ (alternative $A_i$ selon $C_j$). La priorité finale est
\begin{equation}
  P(A_i)=\sum_{j=1}^n w_j\, w_{ji}.
\end{equation}

\paragraph{Implémentation.} La fonction \texttt{ahp(A, method)} propose la méthode approximative du cours
(par défaut) et le calcul exact du vecteur propre par \texttt{numpy.linalg.eig}. L'utilisateur saisit
uniquement le triangle supérieur (les fractions telles que \texttt{1/3} sont acceptées) ; le triangle
inférieur est complété par réciprocité, ce qui rend impossible une matrice non réciproque. Pour $n=1,2$,
$RI=0$ et la matrice est toujours consistante ; au-delà de $n=10$ on utilise les valeurs étendues de Saaty.
La fonction \texttt{ahp\_rank} enchaîne les priorités locales et globales.

\section{BWM (Best-Worst Method)}

\subsection{Principe}
L'expert désigne le critère le plus important $B$ et le moins important $W$, puis ne compare chaque critère
qu'à ces deux repères, sur l'échelle 1 à 9 :
\begin{itemize}
  \item vecteur \emph{Best-to-Others} $a_{B\cdot}=(a_{B1},\dots,a_{Bn})$ ;
  \item vecteur \emph{Others-to-Worst} $a_{\cdot W}=(a_{1W},\dots,a_{nW})$.
\end{itemize}
BWM demande $2n-3$ comparaisons contre $n(n-1)/2$ pour AHP. Deux contrôles de bon sens s'appliquent :
$a_{BB}=a_{WW}=1$ et la note la plus élevée du vecteur BO doit porter sur le critère $W$.

\subsection{Modèle linéaire du cours}
Idéalement $w_B/w_j=a_{Bj}$ et $w_j/w_W=a_{jW}$. Les jugements n'étant jamais parfaitement cohérents, on
minimise le plus grand écart $\xi$ :
\begin{equation}
\begin{aligned}
  \min\ & \xi\\
  \text{s.c. }& |w_B-a_{Bj}\,w_j|\le\xi,\quad |w_j-a_{jW}\,w_W|\le\xi,\quad \forall j,\\
  & \textstyle\sum_j w_j=1,\quad w_j\ge 0,\quad \xi\ge 0.
\end{aligned}
\end{equation}
Chaque valeur absolue est linéarisée en deux contraintes, ce qui donne la forme matricielle
$A\,x\le b$ avec $x=(w_1,\dots,w_n,\xi)$ :
\begin{equation}
\begin{aligned}
  w_B-a_{Bj}w_j-\xi&\le 0, & -w_B+a_{Bj}w_j-\xi&\le 0,\\
  w_j-a_{jW}w_W-\xi&\le 0, & -w_j+a_{jW}w_W-\xi&\le 0.
\end{aligned}
\end{equation}

\paragraph{Implémentation.} Le programme linéaire est résolu par
\texttt{scipy.optimize.linprog} (solveur HiGHS) avec $c=(0,\dots,0,1)$, la contrainte d'égalité
$\sum_j w_j=1$ et les bornes $x\ge 0$. La matrice $A$ et le libellé de chaque contrainte sont affichés dans
l'interface, ce qui reproduit l'exercice \og écrire ces contraintes sous forme matricielle \fg{} du cours.

\subsection{Variante non linéaire (Rezaei, 2015)}
À titre de comparaison, la plateforme propose aussi le modèle original
\begin{equation}
  \min\ \xi\quad\text{s.c.}\quad \left|\frac{w_B}{w_j}-a_{Bj}\right|\le\xi,\quad
  \left|\frac{w_j}{w_W}-a_{jW}\right|\le\xi,\quad \sum_j w_j=1,
\end{equation}
résolu par SLSQP avec plusieurs points de départ, et le ratio de consistance $CR=\xi^*/CI(a_{BW})$.

\section{Méthode de l'entropie}

L'entropie mesure le désordre contenu dans la colonne d'un critère. Une entropie proche de 1 (valeurs
uniformes) signifie que le critère distingue mal les alternatives et reçoit un poids faible.
\begin{equation}
  p_{ij}=\frac{d_{ij}}{\sum_{i=1}^m d_{ij}},\qquad
  E_j=-\frac{1}{\ln m}\sum_{i=1}^m p_{ij}\ln p_{ij},\qquad
  w_j=\frac{1-E_j}{n-\sum_{k=1}^n E_k}.
\end{equation}
La quantité $1-E_j$ est le degré de diversification du critère $j$.

\paragraph{Implémentation.} On adopte la convention $0\ln 0=0$. La méthode exige des valeurs positives ou
nulles ; une erreur explicite est levée sinon. Si toutes les entropies valent 1, on revient à
l'équipondération.

\section{Méthode CRITIC}

CRITIC (\emph{CRiteria Importance Through Intercriteria Correlation}) combine l'intensité du contraste
(dispersion) et le conflit entre critères (corrélation).
\begin{enumerate}
  \item Normalisation min--max selon le sens du critère :
  \begin{equation}
    r_{ij}=\frac{x_{ij}-\min_i x_{ij}}{\max_i x_{ij}-\min_i x_{ij}}\ (j\in J),\qquad
    r_{ij}=\frac{\max_i x_{ij}-x_{ij}}{\max_i x_{ij}-\min_i x_{ij}}\ (j\in J').
  \end{equation}
  \item Coefficient de corrélation de Pearson entre critères :
  \begin{equation}
    \rho_{jk}=\frac{\sum_i (r_{ij}-\bar r_j)(r_{ik}-\bar r_k)}
                   {\sqrt{\sum_i (r_{ij}-\bar r_j)^2\sum_i (r_{ik}-\bar r_k)^2}}.
  \end{equation}
  \item Écart-type et quantité d'information :
  \begin{equation}
    \sigma_j=\sqrt{\frac{1}{m-1}\sum_i (r_{ij}-\bar r_j)^2},\qquad
    C_j=\sigma_j\sum_{k=1}^n\bigl(1-|\rho_{jk}|\bigr).
  \end{equation}
  \item Poids : $w_j=C_j/\sum_k C_k$.
\end{enumerate}
Plus un critère est dispersé et peu corrélé aux autres, plus il est important. Conformément au support, la
valeur absolue $|\rho_{jk}|$ est utilisée et l'écart-type est calculé avec le dénominateur $m-1$.

\section{Pondération hybride}

Deux combinaisons sont proposées entre un vecteur subjectif $w^{s}$ (AHP, BWM ou manuel) et un vecteur
objectif $w^{o}$ (entropie ou CRITIC) :
\begin{equation}
  w_j=\alpha\,w_j^{s}+(1-\alpha)\,w_j^{o}\quad(\alpha\in[0,1]),\qquad\text{ou}\qquad
  w_j=\frac{w_j^{s}\,w_j^{o}}{\sum_k w_k^{s}\,w_k^{o}}.
\end{equation}
La forme linéaire est celle utilisée dans le cours pour la mise à jour des poids par XAI
($w_j^{new}=\alpha w_j^{AHP}+(1-\alpha)|\phi_j|/\sum|\phi_j|$).

% =====================================================================
\chapter{Classement des alternatives}\label{ch:classement}
% =====================================================================

\section{Normalisation linéaire et solutions idéales}

Pour les méthodes de pondération simple, la matrice est normalisée par
\begin{equation}
  r_{ij}=\frac{x_{ij}}{\max_i x_{ij}}\quad (j\in J),\qquad
  r_{ij}=\frac{\min_i x_{ij}}{x_{ij}}\quad (j\in J'),
\end{equation}
de sorte que $r_{ij}\in\,]0,1]$ et que 1 corresponde toujours à la meilleure valeur. Les solutions idéales
positive et négative sont $A^+=\{x_1^+,\dots,x_n^+\}$ et $A^-=\{x_1^-,\dots,x_n^-\}$.

\section{WSM, WPM et WASPAS}

\begin{align}
  \text{WSM :}\quad & Q_i^{(1)}=\sum_{j=1}^n w_j\,r_{ij},\\
  \text{WPM :}\quad & Q_i^{(2)}=\prod_{j=1}^n r_{ij}^{\,w_j},\\
  \text{WASPAS :}\quad & Q_i=\lambda\,Q_i^{(1)}+(1-\lambda)\,Q_i^{(2)},\qquad \lambda\in[0,1].
\end{align}
WASPAS (Zavadskas et al., 2012) correspond à $\lambda=0{,}5$ dans sa forme de base ; la plateforme laisse
$\lambda$ réglable. Pour les trois méthodes, l'alternative de score le plus élevé est la meilleure.

\section{TOPSIS}

TOPSIS (Hwang et Yoon, 1981) choisit l'alternative la plus proche de la solution idéale positive et la plus
éloignée de la solution idéale négative.
\begin{enumerate}
  \item Normalisation vectorielle : $r_{ij}=x_{ij}\big/\sqrt{\sum_i x_{ij}^2}$.
  \item Matrice normalisée pondérée : $v_{ij}=w_j\,r_{ij}$.
  \item Solutions idéales : $v_j^*=\max_i v_{ij}$ si $j\in J$, $\min_i v_{ij}$ si $j\in J'$ ; et
        inversement pour $v_j'$.
  \item Mesures de séparation :
  \begin{equation}
    S_i^+=\sqrt{\sum_j (v_j^*-v_{ij})^2},\qquad S_i^-=\sqrt{\sum_j (v_j'-v_{ij})^2}.
  \end{equation}
  \item Coefficient de proximité : $RC_i^*=\dfrac{S_i^-}{S_i^-+S_i^+}\in[0,1]$ ; le plus élevé est le
        meilleur.
\end{enumerate}

\section{VIKOR}

VIKOR (méthode de classement de compromis) recherche une solution qui maximise l'utilité de groupe et
minimise le regret individuel.
\begin{enumerate}
  \item Meilleures et pires valeurs : $x_j^*=\max_i x_{ij}$, $x_j^-=\min_i x_{ij}$ pour un critère positif
        (inversement pour un critère négatif).
  \item Utilité et regret :
  \begin{equation}
    S_i=\sum_j w_j\frac{x_j^*-x_{ij}}{x_j^*-x_j^-},\qquad
    R_i=\max_j\left[w_j\frac{x_j^*-x_{ij}}{x_j^*-x_j^-}\right].
  \end{equation}
  \item Indice de compromis, avec $S^*=\min S_i$, $S^-=\max S_i$, $R^*=\min R_i$, $R^-=\max R_i$ et
        $v=0{,}5$ par défaut :
  \begin{equation}
    Q_i=v\,\frac{S_i-S^*}{S^--S^*}+(1-v)\,\frac{R_i-R^*}{R^--R^*}.
  \end{equation}
  \item Classement selon $S$, $R$ et $Q$ (plus petit = meilleur).
  \item Solution de compromis $A'$ (plus petit $Q$) si :
  \begin{itemize}
    \item \textbf{C1 --- avantage acceptable} : $Q(A'')-Q(A')\ge DQ=\frac{1}{m-1}$ ;
    \item \textbf{C2 --- stabilité acceptable} : $A'$ est aussi la meilleure selon $S$ et $R$.
  \end{itemize}
  Si seule C2 échoue, le compromis est $\{A',A''\}$ ; si C1 échoue, c'est
  $\{A',\dots,A^{(k)}\}$ avec $Q(A^{(k)})-Q(A')<\frac{1}{m-1}$.
\end{enumerate}

\paragraph{Implémentation.} Lorsque $x_j^*=x_j^-$ (critère constant), le terme correspondant est nul ;
lorsque $S^-=S^*$ ou $R^-=R^*$, la fraction correspondante est nulle. Pour les comparaisons inter-méthodes,
le score \og plus grand = meilleur \fg{} de VIKOR est $1-Q_i$, mais l'interface affiche toujours $Q_i$.

\section{PROMETHEE II (surclassement)}

Le cours place PROMETHEE parmi les méthodes de surclassement, sans compensation entre critères. Les
diapositives correspondantes ne détaillant pas les formules, nous avons implémenté la formulation standard
de Brans et Vincke (1985). Pour chaque paire $(a,b)$ et chaque critère, l'écart orienté
$d_j(a,b)=g_j(a)-g_j(b)$ (signe inversé pour un critère négatif) est transformé en degré de préférence
$P_j(a,b)\in[0,1]$ par l'une des six fonctions généralisées (tableau~\ref{tab:promethee}). On calcule :
\begin{equation}
  \pi(a,b)=\sum_j w_j P_j(a,b),\quad
  \phi^+(a)=\frac{1}{m-1}\sum_{x}\pi(a,x),\quad
  \phi^-(a)=\frac{1}{m-1}\sum_{x}\pi(x,a),\quad
  \phi(a)=\phi^+(a)-\phi^-(a).
\end{equation}
PROMETHEE II classe selon le flux net $\phi$ ; PROMETHEE I fournit un préordre partiel (surclassement,
indifférence, incomparabilité) affiché dans l'interface.

\begin{table}[H]
\centering
\caption{Fonctions de préférence PROMETHEE implémentées}\label{tab:promethee}
\begin{tabular}{lll}
\toprule
\textbf{Type} & \textbf{Nom} & \textbf{$P(d)$ pour $d>0$}\\
\midrule
I & Usuel & $1$\\
II & En U (seuil $q$) & $1$ si $d>q$, $0$ sinon\\
III & En V (seuil $p$) & $\min(d/p,\,1)$\\
IV & Par paliers ($q$, $p$) & $0$ si $d\le q$ ; $1/2$ si $q<d\le p$ ; $1$ sinon\\
V & Linéaire ($q$, $p$) & $0$ si $d\le q$ ; $(d-q)/(p-q)$ si $q<d<p$ ; $1$ sinon\\
VI & Gaussien ($s$) & $1-\exp\left(-d^2/(2s^2)\right)$\\
\bottomrule
\end{tabular}
\end{table}

% =====================================================================
\chapter{Évaluation d'une alternative isolée}\label{ch:single}
% =====================================================================

On évalue une alternative unique $A$ de valeurs $x=(x_1,\dots,x_n)$ et de poids $w$ pour produire un score
$S(A)$ et une décision (accepter / rejeter, classe A--D, niveau de risque).

\section{MAUT}
\begin{equation}
  S(A)=\sum_{j=1}^n w_j\,u_j(x_j),\qquad
  u_j^{+}(x)=\frac{x-x_j^{min}}{x_j^{max}-x_j^{min}},\qquad
  u_j^{-}(x)=\frac{x-x_j^{max}}{x_j^{min}-x_j^{max}}.
\end{equation}
Les bornes $[x_j^{min},x_j^{max}]$ sont saisies par l'utilisateur (ou reprises de la matrice de décision) ;
les utilités sont bornées à $[0,1]$. La décision est \og accepter si $S(A)\ge\tau$, rejeter sinon \fg.

\section{Choix du seuil de décision}

\subsection{Seuil fixé par la politique}
Le seuil $\tau$ est défini par des experts, la réglementation ou les objectifs stratégiques (par exemple
$\tau=0{,}7$). La plateforme propose aussi une classification en classes A ($S\ge 0{,}8$), B ($\ge 0{,}6$),
C ($\ge 0{,}4$) et D.

\subsection{Seuil issu des données : courbe ROC}
À partir d'alternatives passées de scores $S_i$ et de décisions réelles $y_i$, on prédit
$\hat y_i=1$ si $S_i\ge\tau$. Pour chaque seuil :
\begin{equation}
  TPR=\frac{TP}{TP+FN},\qquad FPR=\frac{FP}{FP+TN}.
\end{equation}
Deux critères du cours sont implémentés :
\begin{itemize}
  \item \textbf{distance au point idéal} $(0,1)$ : minimiser $d=\sqrt{FPR^2+(1-TPR)^2}$ ;
  \item \textbf{indice de Youden} : maximiser $J=TPR-FPR$.
\end{itemize}
Nous y ajoutons la maximisation de l'exactitude et la minimisation d'un coût
$C=c_{FP}\cdot FP+c_{FN}\cdot FN$ (\og optimisation sensible aux coûts \fg), qui permet de départager les
seuils à égalité, ainsi que l'aire sous la courbe (AUC, méthode des trapèzes).

\section{Distance à l'idéal}
Même pour une seule alternative, la logique TOPSIS donne un indice de performance :
\begin{equation}
  D^+=\sqrt{\sum_j (x_j-x_j^+)^2},\qquad D^-=\sqrt{\sum_j (x_j-x_j^-)^2},\qquad C=\frac{D^-}{D^++D^-}.
\end{equation}
Avec des critères d'unités différentes, la formule brute est dominée par les critères de grande amplitude.
La plateforme propose donc, en option activée par défaut, de ramener chaque valeur à une utilité
$u_j=(x_j-x_j^-)/(x_j^+-x_j^-)$ et de pondérer les distances : $D^+=\sqrt{\sum_j w_j(1-u_j)^2}$,
$D^-=\sqrt{\sum_j w_j u_j^2}$.

\section{Évaluation floue d'une seule alternative}
Chaque évaluation est un nombre flou triangulaire $\tilde x=(a,b,c)$ ($a$ pessimiste, $b$ le plus probable,
$c$ optimiste) :
\begin{enumerate}
  \item normalisation $\tilde r_j=(a,b,c)/c_j^{max}$ ;
  \item agrégation $\tilde S=\sum_j w_j\,\tilde r_j$ ;
  \item défuzzification par le centre de gravité $S=(a+b+c)/3$.
\end{enumerate}
La largeur $c-a$ de $\tilde S$ mesure l'incertitude et doit être surveillée en plus du score.

% =====================================================================
\chapter{MCDM floue}\label{ch:floue}
% =====================================================================

\section{Rappels sur les nombres flous}

Un nombre flou triangulaire (TFN) $\tilde A=(a_1,a_2,a_3)$ a pour fonction d'appartenance
\begin{equation}
  \mu_{\tilde A}(x)=
  \begin{cases}
    \dfrac{x-a_1}{a_2-a_1} & a_1\le x\le a_2,\\[6pt]
    \dfrac{a_3-x}{a_3-a_2} & a_2\le x\le a_3,\\[4pt]
    0 & \text{sinon.}
  \end{cases}
\end{equation}
Pour deux TFN positifs, $\tilde A\oplus\tilde B=(a_1+b_1,a_2+b_2,a_3+b_3)$,
$\tilde A\otimes\tilde B=(a_1b_1,a_2b_2,a_3b_3)$ et $\lambda\tilde A=(\lambda a_1,\lambda a_2,\lambda a_3)$.

\section{TOPSIS flou (Chen, 2000)}

Les variables linguistiques servent à évaluer à la fois les alternatives et l'importance des critères
(tableau~\ref{tab:chen}).

\begin{table}[H]
\centering
\caption{Échelles linguistiques de Chen (2000)}\label{tab:chen}
\begin{tabular}{llcllc}
\toprule
\multicolumn{3}{c}{\textbf{Évaluation des alternatives}} & \multicolumn{3}{c}{\textbf{Importance des critères}}\\
\cmidrule(lr){1-3}\cmidrule(lr){4-6}
VP & Very poor & (0, 0, 1) & VL & Very low & (0, 0, 0,1)\\
P & Poor & (0, 1, 3) & L & Low & (0, 0,1, 0,3)\\
MP & Medium poor & (1, 3, 5) & ML & Medium low & (0,1, 0,3, 0,5)\\
F & Fair & (3, 5, 7) & M & Medium & (0,3, 0,5, 0,7)\\
MG & Medium good & (5, 7, 9) & MH & Medium high & (0,5, 0,7, 0,9)\\
G & Good & (7, 9, 10) & H & High & (0,7, 0,9, 1)\\
VG & Very good & (9, 10, 10) & VH & Very high & (0,9, 1, 1)\\
\bottomrule
\end{tabular}
\end{table}

\begin{enumerate}
  \item Chaque décideur $k$ évalue chaque alternative selon chaque critère : $\tilde x_{ij}^k$.
  \item Chaque décideur évalue l'importance de chaque critère : $\tilde w_j^k$.
  \item Évaluations collectives (moyenne des $K$ décideurs) :
  \begin{equation}
    \tilde x_{ij}=\frac{1}{K}\left(\tilde x_{ij}^1\oplus\cdots\oplus\tilde x_{ij}^K\right),\qquad
    \tilde w_j=\frac{1}{K}\left(\tilde w_j^1\oplus\cdots\oplus\tilde w_j^K\right).
  \end{equation}
  \item Normalisation :
  \begin{equation}
    \tilde r_{ij}=\left(\frac{a_{ij}}{c_j^*},\frac{b_{ij}}{c_j^*},\frac{c_{ij}}{c_j^*}\right),\ c_j^*=\max_i c_{ij}\ (j\in J);\qquad
    \tilde r_{ij}=\left(\frac{a_j^-}{c_{ij}},\frac{a_j^-}{b_{ij}},\frac{a_j^-}{a_{ij}}\right),\ a_j^-=\min_i a_{ij}\ (j\in J').
  \end{equation}
  \item Matrice pondérée : $\tilde v_{ij}=\tilde r_{ij}\otimes\tilde w_j$.
  \item Solutions idéales : $\tilde v_j^*=(1,1,1)$ et $\tilde v_j^-=(0,0,0)$.
  \item Distances (méthode du sommet) :
  \begin{equation}
    d(\tilde a,\tilde b)=\sqrt{\tfrac13\left[(a_1-b_1)^2+(a_2-b_2)^2+(a_3-b_3)^2\right]},\quad
    D_i^*=\sum_j d(\tilde v_{ij},\tilde v_j^*),\quad D_i^-=\sum_j d(\tilde v_{ij},\tilde v_j^-).
  \end{equation}
  \item Coefficient de proximité : $CC_i=D_i^-/(D_i^-+D_i^*)$ ; on classe par $CC_i$ décroissant.
\end{enumerate}

\begin{remarque}
Pour un critère négatif, la normalisation doit inverser l'ordre des composantes afin que
$\tilde r_{ij}$ reste un TFN ordonné ($r_1\le r_2\le r_3$) : on divise $a_j^-$ par $(c_{ij},b_{ij},a_{ij})$.
C'est l'écriture retenue ci-dessus et dans le code.
\end{remarque}

\section{Ensembles flous intuitionnistes}

Un nombre flou intuitionniste (IFN) $\alpha=(\mu,\nu)$ vérifie $\mu,\nu\in[0,1]$ et $\mu+\nu\le 1$ ;
$\pi=1-\mu-\nu$ est le degré d'hésitation. Les opérations utilisées sont :
\begin{equation}
  \alpha_1\oplus\alpha_2=(\mu_1+\mu_2-\mu_1\mu_2,\ \nu_1\nu_2),\qquad
  \alpha_1\otimes\alpha_2=(\mu_1\mu_2,\ \nu_1+\nu_2-\nu_1\nu_2),
\end{equation}
la fonction score $S(\alpha)=\mu-\nu$, la précision $H(\alpha)=\mu+\nu$, et l'opérateur d'agrégation
\begin{equation}
  IFWA_\sigma(\alpha_1,\dots,\alpha_K)=\left(1-\prod_{k=1}^K(1-\mu_k)^{\sigma_k},\ \prod_{k=1}^K\nu_k^{\sigma_k}\right).
\end{equation}

\section{TOPSIS flou intuitionniste}

\begin{enumerate}
  \item Identifier les alternatives, les critères et les décideurs $D_1,\dots,D_K$ de poids
        $\sigma_k$ ($\sum_k\sigma_k=1$).
  \item Chaque décideur fournit une matrice d'IFN $\alpha_{ij}^k=(\mu_{ij}^k,\nu_{ij}^k)$.
  \item Agrégation des matrices par IFWA : $\alpha_{ij}=IFWA_\sigma(\alpha_{ij}^1,\dots,\alpha_{ij}^K)$.
  \item Agrégation des poids des critères par IFWA : $w_j=IFWA_\sigma(w_j^1,\dots,w_j^K)$.
  \item Matrice pondérée : $\alpha'_{ij}=\alpha_{ij}\otimes w_j$.
  \item Solutions idéales : pour $j\in J$, $I^+=(\max_i\mu'_{ij},\min_i\nu'_{ij})$ et
        $I^-=(\min_i\mu'_{ij},\max_i\nu'_{ij})$ ; inversement pour $j\in J'$.
  \item Mesures de séparation (distance euclidienne normalisée) :
  \begin{equation}
    S_i^{\pm}=\sqrt{\frac{1}{2n}\sum_{j=1}^n\left[(\mu'_{ij}-\mu_j^{\pm})^2+(\nu'_{ij}-\nu_j^{\pm})^2+(\pi'_{ij}-\pi_j^{\pm})^2\right]}.
  \end{equation}
  \item Coefficient de proximité $C_i=S_i^-/(S_i^-+S_i^+)$ et classement décroissant.
\end{enumerate}
Les échelles linguistiques intuitionnistes par défaut sont celles de Boran et al.\ (2009) ; elles sont
modifiables dans l'interface, et la contrainte $\mu+\nu\le 1$ est vérifiée.

% =====================================================================
\chapter{Analyse de robustesse et comparaison des méthodes}\label{ch:robustesse}
% =====================================================================

\section{Corrélation de rang de Spearman}
Pour deux classements des mêmes $n$ alternatives :
\begin{equation}
  \rho=1-\frac{6\sum_i d_i^2}{n(n^2-1)},\qquad d_i=\text{rang}_1(A_i)-\text{rang}_2(A_i).
\end{equation}
$\rho=1$ : classements identiques ; $\rho=0$ : aucun lien ; $\rho=-1$ : classements inversés. En présence
d'ex aequo, la formule n'est plus exacte : la plateforme bascule alors sur
\texttt{scipy.stats.spearmanr}. Le seuil de significativité bilatéral à 5\,\% est affiché pour
$n\le 20$ (par exemple $0{,}886$ pour $n=6$).

\section{Perturbation d'un poids}
Le poids du critère $j$ parcourt $[0,1]$ ; les autres poids sont redistribués proportionnellement à leurs
valeurs initiales pour conserver leurs rapports :
\begin{equation}
  w'_j=v,\qquad w'_k=w_k\,\frac{1-v}{1-w_j}\quad(k\ne j).
\end{equation}
La plateforme trace les scores (ou les rangs) en fonction de $v$ et signale les valeurs de $v$ pour
lesquelles l'alternative classée première change.

\section{Simulation de Monte Carlo}
Des vecteurs de poids $w'\sim\text{Dirichlet}(\kappa\,w)$ sont tirés autour des poids actifs ($\kappa$
règle l'intensité de la perturbation). Pour chaque tirage on recalcule le classement ; on obtient la
fréquence de chaque rang par alternative et la distribution de $\rho$ entre le classement perturbé et le
classement de référence.

\section{Consensus}
Un classement de consensus est obtenu par comptage de Borda : on somme les rangs attribués par chaque
méthode et on classe par somme croissante.

% =====================================================================
\chapter{Architecture et implémentation}\label{ch:architecture}
% =====================================================================

\section{Pile technologique}

\begin{table}[H]
\centering
\caption{Pile technologique}\label{tab:stack}
\begin{tabularx}{\linewidth}{lX}
\toprule
\textbf{Composant} & \textbf{Rôle}\\
\midrule
Python 3.10+ & Langage (testé sous Python 3.13)\\
Streamlit & Interface web : \texttt{st.data\_editor}, \texttt{st.tabs}, \texttt{st.sidebar}, \texttt{st.metric}, \texttt{st.latex}\\
NumPy & Calcul matriciel (normalisations, distances, valeurs propres)\\
SciPy & \texttt{optimize.linprog} (BWM linéaire), \texttt{optimize.minimize} (BWM non linéaire), \texttt{stats.rankdata}, \texttt{stats.spearmanr}\\
Pandas & Tableaux affichés et édités\\
Plotly & Graphiques interactifs : radar, barres, courbes de sensibilité, courbe ROC, cartes de chaleur\\
\bottomrule
\end{tabularx}
\end{table}

\section{Structure du projet}

\begin{lstlisting}[language={},numbers=none,caption={Arborescence du projet}]
mcdm_app/
|-- app.py                     # point d'entree, navigation, etat global
|-- requirements.txt
|-- core/                      # moteur de calcul (aucune dependance a Streamlit)
|   |-- weighting.py           # AHP, BWM (lineaire et non lineaire), Entropie, CRITIC, hybride
|   |-- crisp_ranking.py       # normalisations, WSM, WPM, WASPAS, TOPSIS, VIKOR
|   |-- outranking.py          # PROMETHEE I / II
|   |-- single_alternative.py  # MAUT, seuil, ROC / Youden, distance a l'ideal, TFN
|   |-- fuzzy_mcdm.py          # TOPSIS flou (Chen 2000), IFWA, TOPSIS flou intuitionniste
|   `-- sensitivity.py         # perturbation, Monte Carlo, Spearman, Borda
|-- utils/
|   |-- presets.py             # exercices du cours pre-charges
|   `-- validation.py          # controles de dimension, de coherence, de domaine
|-- ui/                        # une page Streamlit par module
|   |-- common.py              # etat de session, editeurs persistants, graphiques
|   |-- page_space.py  page_weighting.py  page_evaluation.py
|   `-- page_fuzzy.py  page_sensitivity.py  page_cases.py
|-- tests/
|   `-- test_course_examples.py
`-- rapport/
    `-- rapport.tex
\end{lstlisting}

\begin{figure}[H]
\centering
\begin{tikzpicture}[
  node distance=0.5cm and 0.6cm,
  layer/.style={draw=emiblue, rounded corners=3pt, align=center, font=\small, minimum height=0.9cm},
  ui/.style={layer, fill=lightblue!60, text width=2.3cm},
  core/.style={layer, fill=white, text width=2.3cm},
  arr/.style={-{Stealth[length=2mm]}, emiblue}]
  \node[layer, fill=emiblue, text=white, text width=12.2cm] (app) {\texttt{app.py} : navigation (barre latérale) et état de session};
  \node[ui, below=0.6cm of app.south west, anchor=north west] (u1) {Espace de décision};
  \node[ui, right=0.25cm of u1] (u2) {Pondération};
  \node[ui, right=0.25cm of u2] (u3) {Évaluation};
  \node[ui, right=0.25cm of u3] (u4) {Floue / Sensibilité / Cas};
  \node[core, below=0.9cm of u1] (c1) {\texttt{validation}\\\texttt{presets}};
  \node[core, below=0.9cm of u2] (c2) {\texttt{weighting}};
  \node[core, below=0.9cm of u3] (c3) {\texttt{crisp\_ranking}\\\texttt{outranking}\\\texttt{single\_alternative}};
  \node[core, below=0.9cm of u4] (c4) {\texttt{fuzzy\_mcdm}\\\texttt{sensitivity}};
  \foreach \i in {1,2,3,4} { \draw[arr] (app.south -| u\i.north) -- (u\i.north); \draw[arr] (u\i) -- (c\i); }
  \node[draw=codecomment, dashed, rounded corners, fit=(c1)(c4), inner sep=6pt, label={[font=\small\itshape]below:couche \texttt{core} : NumPy / SciPy, testable sans interface}] {};
\end{tikzpicture}
\caption{Architecture en couches de la plateforme}\label{fig:archi}
\end{figure}

\section{Principes de conception}

\begin{itemize}
  \item \textbf{Séparation calcul / interface.} Le paquet \texttt{core} ne dépend que de NumPy, SciPy et
        Pandas. Chaque fonction reçoit des tableaux et renvoie un dictionnaire contenant le résultat
        \emph{et} toutes les grandeurs intermédiaires (matrices normalisées, solutions idéales,
        distances\dots), ce qui permet à l'interface d'afficher chaque étape du cours et aux tests de
        vérifier les calculs sans lancer Streamlit.
  \item \textbf{Conventions uniformes.} Les types de critères sont codés \texttt{"+"} / \texttt{"-"} ;
        toutes les méthodes renvoient une clé \texttt{scores} (plus grand = meilleur) et une clé
        \texttt{ranks} (1 = meilleure, ex aequo au rang minimal via \texttt{scipy.stats.rankdata}). Cette
        interface commune permet aux modules de sensibilité et de consensus d'appeler n'importe quelle
        méthode de façon générique.
  \item \textbf{Fidélité au cours.} Chaque formule est commentée avec la notation du support ; lorsque
        nous nous en écartons (normalisation floue des critères négatifs, option de normalisation pour la
        distance à l'idéal), l'écart est signalé dans l'interface et dans ce rapport.
  \item \textbf{Robustesse numérique.} Les divisions par zéro (critère constant, colonne nulle) sont
        traitées explicitement ; les domaines de validité sont contrôlés (valeurs strictement positives
        pour WPM et la normalisation $\min/x$, $a\le b\le c$ pour les TFN, $\mu+\nu\le1$ pour les IFN,
        réciprocité et échelle de Saaty pour AHP, contrôles de bon sens BWM).
\end{itemize}

\section{Gestion de l'état et des tableaux éditables}

Streamlit réexécute le script à chaque interaction et détruit l'état d'un widget dès qu'il n'est plus
affiché (changement de page). Pour que les tableaux saisis par l'utilisateur survivent à la navigation, nous
avons écrit un \emph{éditeur persistant} : la valeur éditée est recopiée dans \texttt{st.session\_state},
et la base de l'éditeur n'est réinitialisée à partir de cette copie que lorsque le widget a disparu. Un
numéro de version incrémenté à chaque changement de structure (chargement d'un exercice, ajout d'un
critère) force la recréation propre de l'éditeur.

\begin{lstlisting}[caption={Editeur persistant (ui/common.py)}]
def persistent_editor(name, default, **kwargs):
    store, base = f"_val_{name}", f"_base_{name}"
    key = f"_ed_{name}_{ss.get(f'_v_{name}', 0)}"
    if store not in ss:
        ss[store] = default.copy()
    if key not in ss:                 # widget absent : on repart de la derniere valeur
        ss[base] = ss[store].copy()
    edited = st.data_editor(ss[base], key=key, **kwargs)
    ss[store] = edited
    return edited
\end{lstlisting}

Lorsque l'utilisateur renomme, ajoute ou supprime un critère ou une alternative, la matrice de décision et
le vecteur de poids sont réindexés : les valeurs sont conservées par nom, ou par position en cas de
renommage, et les nouveaux critères reçoivent le poids moyen avant renormalisation.

\section{Extraits de code}

\begin{lstlisting}[caption={BWM lineaire : construction de A x <= b et resolution (core/weighting.py)}]
def bwm(best, worst, a_BO, a_OW):
    n = len(a_BO)
    c = np.zeros(n + 1); c[-1] = 1.0          # min xi
    rows = []
    for j in range(n):
        r = np.zeros(n + 1); r[best] += 1; r[j] -= a_BO[j]; r[-1] = -1; rows.append(r)
        r = np.zeros(n + 1); r[best] -= 1; r[j] += a_BO[j]; r[-1] = -1; rows.append(r)
        r = np.zeros(n + 1); r[j] += 1; r[worst] -= a_OW[j]; r[-1] = -1; rows.append(r)
        r = np.zeros(n + 1); r[j] -= 1; r[worst] += a_OW[j]; r[-1] = -1; rows.append(r)
    A_eq = np.zeros((1, n + 1)); A_eq[0, :n] = 1.0   # somme des poids = 1
    res = linprog(c, A_ub=np.array(rows), b_ub=np.zeros(len(rows)),
                  A_eq=A_eq, b_eq=[1.0], bounds=[(0, None)] * (n + 1), method="highs")
    return res.x[:n], res.x[-1]
\end{lstlisting}

\begin{lstlisting}[caption={TOPSIS (core/crisp\_ranking.py, version simplifiee)}]
def topsis(X, w, types):
    R = X / np.sqrt((X ** 2).sum(axis=0))                  # etape 1
    V = R * w                                              # etape 2
    b = np.array([t == "+" for t in types])
    v_plus  = np.where(b, V.max(axis=0), V.min(axis=0))    # etape 3
    v_minus = np.where(b, V.min(axis=0), V.max(axis=0))
    s_plus  = np.sqrt(((V - v_plus) ** 2).sum(axis=1))     # etape 4
    s_minus = np.sqrt(((V - v_minus) ** 2).sum(axis=1))
    return s_minus / (s_plus + s_minus)                    # etape 5
\end{lstlisting}

% =====================================================================
\chapter{Interface utilisateur}\label{ch:interface}
% =====================================================================

L'application se lance par \texttt{streamlit run app.py}. Une barre latérale donne accès à sept pages,
rappelle le problème courant (dimensions, source des poids actifs) et permet de charger un exercice du cours
en un clic.

\begin{description}
  \item[Accueil] Présentation du processus, de la taxonomie et du guide de choix des méthodes.
  \item[1 · Espace de décision] Deux tableaux éditables (critères avec leur type $+/-$, alternatives) et la
        matrice de décision. Import et export CSV (une ligne \texttt{type} optionnelle contient les sens de
        préférence). Contrôles automatiques : valeurs manquantes, noms dupliqués, critère constant, valeurs
        non positives. Un graphique radar affiche les profils normalisés min--max des alternatives.
  \item[2 · Pondération] Un onglet par méthode : manuelle, AHP (saisie du triangle supérieur, méthode
        approximative ou exacte, $\lambda_{max}$, $CI$, $RI$, $CR$ et verdict), BWM (choix de Best et Worst,
        vecteurs BO et OW, contrôles de bon sens, modèle linéaire ou non linéaire, $\xi^*$, forme matricielle
        $Ax\le b$), Entropie, CRITIC (avec carte de chaleur des corrélations) et hybride. Le bouton
        \og Appliquer ces poids \fg{} rend un vecteur actif pour toutes les pages ; un dernier onglet
        compare graphiquement tous les vecteurs calculés.
  \item[3 · Évaluation] Choix de la portée. \emph{Plusieurs alternatives} : sélection des méthodes (WSM,
        WPM, WASPAS, TOPSIS, VIKOR, PROMETHEE II), paramètres ($\lambda$, $v$, fonctions de préférence par
        critère), tableau de synthèse des rangs et un onglet détaillé par méthode (formules en
        \LaTeX, tableaux intermédiaires, conditions C1/C2 de VIKOR, préordre PROMETHEE I, graphique des
        scores). \emph{Une seule alternative} : MAUT avec seuil et classes, seuil par ROC (tableau des seuils
        testés, optimum de Youden, distance, exactitude, coût, AUC, courbe ROC), distance à l'idéal et
        évaluation floue.
  \item[4 · MCDM floue] TOPSIS flou de Chen et TOPSIS flou intuitionniste : nombre de décideurs réglable, un
        onglet de saisie linguistique par décideur (listes déroulantes), échelles modifiables, poids des
        décideurs, et résolution étape par étape conforme aux questions du cours.
  \item[5 · Sensibilité \& consensus] Perturbation d'un poids (scores ou rangs en fonction du poids,
        détection des inversions), Monte Carlo (carte des fréquences de rang, histogramme de $\rho$,
        probabilité que la meilleure alternative le reste), consensus (tableau des rangs, Borda, matrice de
        Spearman, significativité).
  \item[6 · Études de cas] Dix exercices du cours résolus pas à pas, avec comparaison automatique aux
        valeurs publiées.
\end{description}

Les graphiques utilisent une palette catégorielle à ordre fixe (une couleur par alternative, identique
d'un graphique à l'autre), une rampe séquentielle bleue pour les fréquences et une rampe divergente
rouge--gris--bleu centrée sur 0 pour les corrélations. Tous sont interactifs (info-bulles au survol).

% =====================================================================
\chapter{Validation sur les études de cas du cours}\label{ch:validation}
% =====================================================================

Chaque exercice des diapositives a été saisi dans \texttt{utils/presets.py} et résolu par la plateforme.
Les tableaux ci-dessous reproduisent les résultats calculés (arrondis à trois ou quatre décimales).

\section{AHP --- achat d'une voiture}

Critères : coût, confort, sécurité. Matrice de comparaison :
$\begin{pmatrix}1&7&3\\1/7&1&1/3\\1/3&3&1\end{pmatrix}$, sommes des colonnes $(1{,}476;\ 11;\ 4{,}333)$.

\begin{table}[H]
\centering
\caption{AHP : matrice normalisée et poids des critères}\label{tab:ahp}
\begin{tabular}{lcccc}
\toprule
 & Coût & Confort & Sécurité & \textbf{Poids $w$}\\
\midrule
Coût & 0,6774 & 0,6364 & 0,6923 & \textbf{0,6687}\\
Confort & 0,0968 & 0,0909 & 0,0769 & \textbf{0,0882}\\
Sécurité & 0,2258 & 0,2727 & 0,2308 & \textbf{0,2431}\\
\bottomrule
\end{tabular}
\end{table}

$\lambda_{max}=3{,}0108$, $CI=0{,}0054$, $RI=0{,}58$, $CR=0{,}0093<0{,}1$ : la matrice est consistante.
Le calcul exact du vecteur propre donne $(0{,}6694;\ 0{,}0879;\ 0{,}2426)$ et $\lambda_{max}=3{,}0070$,
très proche de l'approximation.

\begin{table}[H]
\centering
\caption{AHP : priorités locales et finales des voitures}\label{tab:ahp2}
\begin{tabular}{lcccc}
\toprule
 & Coût (0,669) & Confort (0,088) & Sécurité (0,243) & \textbf{Priorité finale}\\
\midrule
Voiture 1 & 0,875 & 0,167 & 0,100 & \textbf{0,624}\\
Voiture 2 & 0,125 & 0,833 & 0,900 & 0,376\\
\bottomrule
\end{tabular}
\end{table}

La voiture 1 est retenue : le coût, critère dominant, l'emporte sur le confort et la sécurité. Les matrices
locales $2\times 2$ sont trivialement consistantes ($CR=0$).

\section{BWM --- choix d'un fournisseur}

Best $=C_1$ (qualité), Worst $=C_4$ (RSE), $a_{B\cdot}=(1,3,4,8)$, $a_{\cdot W}=(8,5,3,1)$.

\begin{table}[H]
\centering
\caption{BWM : poids obtenus et comparaison avec le support}\label{tab:bwm}
\begin{tabular}{lccccc}
\toprule
 & Qualité & Prix & Délai & RSE & $\xi^*$\\
\midrule
Modèle linéaire du cours (plateforme) & 0,563 & 0,216 & 0,162 & 0,060 & 0,0838\\
Modèle non linéaire (Rezaei 2015) & 0,536 & 0,251 & 0,153 & 0,061 & 0,860\\
Valeurs de la diapositive (p.~42) & 0,552 & 0,207 & 0,155 & 0,086 & 0,08\\
\bottomrule
\end{tabular}
\end{table}

La valeur optimale $\xi^*=0{,}084\approx0{,}08$ du cours est retrouvée, de même que l'ordre des critères
(qualité $>$ prix $>$ délai $>$ RSE) et le constat que la qualité pèse plus de la moitié de la décision. En
revanche, les poids affichés sur la diapositive ne sont pas une solution du modèle présenté : avec
$w=(0{,}552;0{,}207;0{,}155;0{,}086)$, on a par exemple $|w_1-8w_4|=0{,}136$ et $|w_2-5w_4|=0{,}223$,
ce qui exigerait $\xi\ge0{,}223$. Nous avons vérifié que l'optimum du programme linéaire est \emph{unique}
(en minimisant puis maximisant chaque $w_j$ sous la contrainte $\xi=\xi^*$, on obtient la même valeur), et
que le modèle non linéaire ne redonne pas non plus ces poids. La plateforme retient donc la solution exacte
du modèle du cours et signale l'écart dans l'étude de cas.

\section{CRITIC --- machines de découpage}

Types : $C_1(-)$, $C_2(-)$, $C_3(+)$, $C_4(-)$, tels qu'indiqués dans l'énoncé.

\begin{table}[H]
\centering
\caption{CRITIC : grandeurs intermédiaires et poids}\label{tab:critic}
\begin{tabular}{lcccc}
\toprule
 & $C_1$ & $C_2$ & $C_3$ & $C_4$\\
\midrule
$\sigma_j$ & 0,4576 & 0,4178 & 0,4787 & 0,4590\\
$\sum_k(1-|\rho_{jk}|)$ & 0,6876 & 0,8368 & 1,2082 & 2,0680\\
$C_j$ & 0,3146 & 0,3496 & 0,5784 & 0,9491\\
\textbf{Poids $w_j$} & \textbf{0,1436} & \textbf{0,1595} & \textbf{0,2639} & \textbf{0,4331}\\
\midrule
Entropie (comparaison) & 0,4109 & 0,3337 & 0,1258 & 0,1297\\
\bottomrule
\end{tabular}
\end{table}

$C_1$ et $C_2$ sont fortement corrélés ($\rho_{12}=0{,}91$) et apportent donc une information redondante ;
$C_4$, presque indépendant des autres, reçoit le poids le plus élevé. La comparaison avec l'entropie
illustre le constat du cours : deux méthodes objectives peuvent donner des poids très différents, car
l'entropie ne mesure que la dispersion relative des valeurs brutes (très forte pour $C_1$, de 30 à 300) sans
tenir compte du sens des critères ni des corrélations. En prolongement, TOPSIS avec les poids CRITIC classe
$A_3$ en tête ($RC^*=0{,}745$), devant $A_1$ (0,599), $A_4$ (0,354) et $A_2$ (0,319).

\section{WSM, WPM et WASPAS --- choix d'une machine}

Types $(-,-,+,-,+)$ et poids $(0{,}331;\ 0{,}181;\ 0{,}369;\ 0{,}072;\ 0{,}047)$.

\begin{table}[H]
\centering
\caption{WSM, WPM et WASPAS ($\lambda=0{,}5$)}\label{tab:waspas}
\begin{tabular}{lcccc}
\toprule
 & $Q^{(1)}$ WSM & $Q^{(2)}$ WPM & $Q$ WASPAS & Rang\\
\midrule
$A_1$ & 0,7443 & 0,7289 & 0,7366 & 4\\
$A_2$ & 0,8310 & 0,8037 & 0,8174 & 3\\
$A_3$ & 0,8746 & 0,8630 & 0,8688 & 2\\
$A_4$ & \textbf{0,9110} & \textbf{0,9060} & \textbf{0,9085} & \textbf{1}\\
\bottomrule
\end{tabular}
\end{table}

Les trois méthodes concordent : la machine $A_4$ est la meilleure.

\section{TOPSIS --- choix d'une voiture}

Critères : style (+), fiabilité (+), économie de carburant (+), coût ($-$) ; poids
$(0{,}2;\ 0{,}1;\ 0{,}4;\ 0{,}3)$.

\begin{table}[H]
\centering
\caption{TOPSIS : matrice normalisée pondérée et solutions idéales}\label{tab:topsis1}
\begin{tabular}{lcccc}
\toprule
 & Style & Fiabilité & Économie & Coût\\
\midrule
M1 & 0,0923 & 0,0614 & 0,2179 & 0,1583\\
M2 & 0,1055 & 0,0477 & 0,1937 & 0,1385\\
M3 & 0,1187 & 0,0409 & 0,1937 & 0,1780\\
M4 & 0,0791 & 0,0477 & 0,1937 & 0,1187\\
\midrule
$I^+$ & 0,1187 & 0,0614 & 0,2179 & 0,1187\\
$I^-$ & 0,0791 & 0,0409 & 0,1937 & 0,1780\\
\bottomrule
\end{tabular}
\end{table}

\begin{table}[H]
\centering
\caption{TOPSIS : séparations et coefficient de proximité}\label{tab:topsis2}
\begin{tabular}{lcccc}
\toprule
 & $S^+$ & $S^-$ & $RC^*$ & Rang\\
\midrule
M1 & 0,0475 & 0,0396 & 0,4545 & 3\\
M2 & 0,0366 & 0,0480 & \textbf{0,5678} & \textbf{1}\\
M3 & 0,0673 & 0,0396 & 0,3703 & 4\\
M4 & 0,0483 & 0,0597 & 0,5527 & 2\\
\bottomrule
\end{tabular}
\end{table}

La marque M2 est sélectionnée, de peu devant M4.

\section{VIKOR --- achat d'une machine}

$C_1(-)$, $C_2(+)$, poids égaux. $x^*=(1;\ 4500)$, $x^-=(5;\ 3000)$.

\begin{table}[H]
\centering
\caption{VIKOR : utilité, regret et indice de compromis}\label{tab:vikor}
\begin{tabular}{lccccc}
\toprule
 & $C_1$ & $C_2$ & $S$ & $R$ & $Q$\\
\midrule
$A_1$ & 1 & 3000 & 0,500 & 0,500 & 1\\
$A_2$ & 2 & 3750 & \textbf{0,375} & \textbf{0,250} & \textbf{0}\\
$A_3$ & 5 & 4500 & 0,500 & 0,500 & 1\\
\bottomrule
\end{tabular}
\end{table}

$A_2$ minimise $Q$. Avantage acceptable : $Q(A'')-Q(A')=1\ge DQ=1/(m-1)=0{,}5$ ; stabilité acceptable :
$A_2$ est aussi la meilleure selon $S$ et $R$. La solution de compromis est $A_2$ seule.

\section{Corrélation de Spearman --- comparaison de méthodes}

\begin{table}[H]
\centering
\caption{Matrice de Spearman (rangs du cours, p.~78)}\label{tab:spearman}
\begin{tabular}{lccc}
\toprule
 & TOPSIS & VIKOR & PROMETHEE II\\
\midrule
TOPSIS & 1,00 & 0,77 & 0,94\\
VIKOR & 0,77 & 1,00 & 0,60\\
PROMETHEE II & 0,94 & 0,60 & 1,00\\
\bottomrule
\end{tabular}
\end{table}

Les trois coefficients du cours sont retrouvés exactement (par exemple TOPSIS--VIKOR :
$\sum d_i^2=8$, $\rho=1-48/210=0{,}771$). Avec $n=6$, seul $\rho\ge0{,}886$ est significatif à 5\,\% :
uniquement la paire TOPSIS--PROMETHEE II.

\section{Seuil de décision --- courbe ROC}

\begin{table}[H]
\centering
\caption{Seuils testés sur les six décisions passées}\label{tab:roc}
\begin{tabular}{lccccc}
\toprule
 & Seuil $\tau$ & TPR & FPR & $d$ & $J$\\
\midrule
Cas 1 & 0,75 & 0,67 & 0,00 & 0,33 & 0,67\\
Cas 2 & 0,65 & 1,00 & 0,33 & 0,33 & 0,67\\
Cas 3 & 0,60 & 1,00 & 0,67 & 0,67 & 0,33\\
\bottomrule
\end{tabular}
\end{table}

Le tableau du cours est reproduit à l'identique ; les cas 1 et 2 sont à égalité et se départagent par le
coût des erreurs. En balayant tous les seuils candidats (les scores observés), la plateforme met en
évidence que $\tau=0{,}71$ sépare parfaitement les deux classes ($TPR=1$, $FPR=0$, $J=1$, $d=0$,
$AUC=1$) : le choix parmi les seuils du cours n'est donc qu'une illustration, la plage optimale se situant
dans $]0{,}69;\ 0{,}71]$.

\section{Évaluation floue d'une seule alternative}

ROA $=(1{,}0;\ 1{,}2;\ 1{,}5)$, liquidité $=(110;\ 120;\ 130)$, poids $(0{,}6;\ 0{,}4)$. Après
normalisation : $(0{,}67;\ 0{,}80;\ 1{,}00)$ et $(0{,}85;\ 0{,}92;\ 1{,}00)$. La plateforme obtient
$\tilde S=(0{,}74;\ 0{,}85;\ 1{,}00)$, $S=0{,}86$ et une largeur $c-a=0{,}26$, valeurs identiques à
celles du cours.

\section{TOPSIS flou (Chen 2000) --- recrutement d'un ingénieur}

Trois candidats, trois décideurs, cinq critères (tous positifs).

\begin{table}[H]
\centering
\caption{TOPSIS flou : poids collectifs des critères}\label{tab:chenw}
\resizebox{\linewidth}{!}{%
\begin{tabular}{lccccc}
\toprule
 & $C_1$ & $C_2$ & $C_3$ & $C_4$ & $C_5$\\
\midrule
$\tilde w_j$ & (0,70; 0,87; 0,97) & (0,90; 1; 1) & (0,77; 0,93; 1) & (0,90; 1; 1) & (0,43; 0,63; 0,83)\\
\bottomrule
\end{tabular}}
\end{table}

\begin{table}[H]
\centering
\caption{TOPSIS flou : matrice de décision collective}\label{tab:chenx}
\resizebox{\linewidth}{!}{%
\begin{tabular}{lccccc}
\toprule
 & $C_1$ & $C_2$ & $C_3$ & $C_4$ & $C_5$\\
\midrule
$A_1$ & (5,67; 7,67; 9,33) & (5,00; 7,00; 8,67) & (5,67; 7,67; 9,00) & (8,33; 9,67; 10) & (3,00; 5,00; 7,00)\\
$A_2$ & (6,33; 8,33; 9,67) & (9,00; 10; 10) & (8,33; 9,67; 10) & (9,00; 10; 10) & (7,00; 8,67; 9,67)\\
$A_3$ & (6,33; 8,00; 9,00) & (7,00; 8,67; 9,67) & (7,00; 8,67; 9,67) & (7,00; 8,67; 9,67) & (6,33; 8,33; 9,67)\\
\bottomrule
\end{tabular}}
\end{table}

Tous les critères étant positifs, la normalisation divise chaque TFN par $c_j^*=\max_i c_{ij}$, soit
$c_j^*=9{,}67$ pour $C_1$ et $C_5$ et $c_j^*=10$ pour $C_2$, $C_3$ et $C_4$ ; la matrice normalisée est
ensuite pondérée par $\tilde w_j$.

\begin{table}[H]
\centering
\caption{TOPSIS flou : distances et coefficients de proximité}\label{tab:chencc}
\begin{tabular}{lcccc}
\toprule
 & $D^*$ & $D^-$ & $CC$ & Rang\\
\midrule
$A_1$ & 1,946 & 3,430 & 0,638 & 3\\
$A_2$ & 1,259 & 4,106 & \textbf{0,765} & \textbf{1}\\
$A_3$ & 1,602 & 3,771 & 0,702 & 2\\
\bottomrule
\end{tabular}
\end{table}

Le classement $A_2\succ A_3\succ A_1$ désigne le candidat $A_2$, conformément à l'exemple de référence de
Chen (2000).

\section{TOPSIS flou intuitionniste (exemple illustratif)}

Le cours ne fournissant pas de données chiffrées pour cette méthode, nous avons construit un exemple de
sélection de fournisseur (4 fournisseurs, 3 décideurs de poids $0{,}40$, $0{,}35$ et $0{,}25$, critères
qualité, délai, service (+) et prix ($-$)). Les poids agrégés des critères par IFWA sont
$(0{,}874;0{,}119)$, $(0{,}681;0{,}266)$, $(0{,}608;0{,}339)$ et $(0{,}856;0{,}132)$. Les coefficients
de proximité obtenus sont $C=(0{,}794;\ 0{,}708;\ 0{,}407;\ 0{,}114)$, d'où $F_1\succ F_2\succ F_3\succ
F_4$. Les tests unitaires vérifient en outre des propriétés de l'opérateur (idempotence de IFWA).

\section{Synthèse}

\begin{table}[H]
\centering
\caption{Synthèse de la validation}\label{tab:synthese}
\begin{tabularx}{\linewidth}{lXl}
\toprule
\textbf{Cas (page)} & \textbf{Grandeur comparée} & \textbf{Résultat}\\
\midrule
AHP voiture (32--34) & $CR<0{,}1$, voiture 1 retenue & \ok\\
BWM fournisseur (37--43) & $\xi^*\approx0{,}08$ et ordre des critères & \ok\\
 & poids numériques de la p.~42 & \ecart{} (commenté)\\
CRITIC (51) & poids, $C_4$ dominant & calculé\\
WASPAS (58) & $A_4$ meilleure pour WSM, WPM, WASPAS & calculé\\
TOPSIS (65) & M2 sélectionnée & calculé\\
VIKOR (72) & $A_2$ compromis, C1 et C2 vérifiées & calculé\\
Spearman (78--79) & 0,77 ; 0,94 ; 0,60 & \ok\\
ROC (85--88) & TPR, FPR, $d$, $J$ des trois cas & \ok\\
Floue une alternative (134) & $\tilde S=(0{,}74;0{,}85;1)$, $S=0{,}86$ & \ok\\
Chen (140--142) & $A_2\succ A_3\succ A_1$ & calculé\\
\bottomrule
\end{tabularx}
\end{table}

\noindent\emph{Calculé} signifie que le cours pose l'exercice sans en publier la solution ; les résultats ont
alors été vérifiés par des propriétés indépendantes (somme des poids, flux nets nuls, cohérence entre
méthodes) et par un recalcul manuel des premières étapes.

% =====================================================================
\chapter{Analyse comparative et de sensibilité}\label{ch:comparaison}
% =====================================================================

Pour illustrer les bonnes pratiques du cours, nous appliquons les six méthodes de classement au problème
TOPSIS du choix d'une voiture (poids $(0{,}2;0{,}1;0{,}4;0{,}3)$), puis étudions la robustesse du
classement.

\section{Comparaison des méthodes}

\begin{table}[H]
\centering
\caption{Rangs obtenus par les six méthodes (choix d'une voiture)}\label{tab:comparaison}
\begin{tabular}{lcccccc|c}
\toprule
 & WSM & WPM & WASPAS & TOPSIS & VIKOR & PROMETHEE II & Borda\\
\midrule
M1 & 1 & 1 & 1 & 3 & 1 & 1 & \textbf{1}\\
M2 & 2 & 2 & 2 & 1 & 2 & 2 & 2\\
M3 & 4 & 4 & 4 & 4 & 4 & 4 & 4\\
M4 & 3 & 3 & 3 & 2 & 3 & 3 & 3\\
\bottomrule
\end{tabular}
\end{table}

Cinq méthodes sur six désignent M1 ; seule TOPSIS place M2 en tête. Le coefficient de Spearman entre TOPSIS
et chacune des autres méthodes vaut $\rho=0{,}40$ ($\sum d_i^2=6$), alors que les cinq autres méthodes
sont parfaitement concordantes entre elles ($\rho=1$). Pour VIKOR, $Q=(0;\ 0{,}821;\ 1;\ 0{,}857)$ :
M1 est solution de compromis et vérifie C1 ($0{,}821\ge 1/3$) et C2. Pour PROMETHEE II (fonction
usuelle), les flux nets sont $\phi=(0{,}333;\ 0{,}033;\ -0{,}333;\ -0{,}033)$.

\paragraph{Interprétation.} La divergence de TOPSIS s'explique par deux différences de construction :
la normalisation vectorielle (au lieu de $x/\max$ ou $\min/x$) et l'agrégation par distances euclidiennes
aux deux solutions idéales. M1 domine fortement sur l'économie de carburant (critère le plus lourd) et la
fiabilité, mais elle est pénalisée sur le coût, où elle se rapproche de l'anti-idéal ; la distance
euclidienne amplifie cet écart, alors que les méthodes additives le compensent. C'est exactement la
situation que le cours décrit : un même problème peut produire des classements différents selon la
méthode, d'où la nécessité de les comparer et de justifier le choix final.

\section{Sensibilité aux poids (TOPSIS)}

\begin{table}[H]
\centering
\caption{Alternative classée première selon le poids d'un critère (TOPSIS, pas de 0,005)}\label{tab:sens}
\begin{tabular}{lcl}
\toprule
\textbf{Critère} & \textbf{Poids actuel} & \textbf{Alternative en tête selon l'intervalle du poids}\\
\midrule
Style & 0,20 & M4 sur $[0;\ 0{,}185]$, M2 sur $[0{,}19;\ 0{,}415]$, M3 sur $[0{,}42;\ 1]$\\
Fiabilité & 0,10 & M2 sur $[0;\ 0{,}18]$, M1 sur $[0{,}185;\ 1]$\\
Économie carburant & 0,40 & M2 sur $[0;\ 0{,}47]$, M4 sur $[0{,}475;\ 0{,}515]$, M1 sur $[0{,}52;\ 1]$\\
Coût & 0,30 & M1 sur $[0;\ 0{,}185]$, M2 sur $[0{,}19;\ 0{,}32]$, M4 sur $[0{,}325;\ 1]$\\
\bottomrule
\end{tabular}
\end{table}

Le poids actuel du style ($0{,}20$) et celui du coût ($0{,}30$) sont tous deux à moins de $0{,}03$ d'un point
d'inversion : la sélection de M2 par TOPSIS est donc \textbf{fragile}.

\section{Simulation de Monte Carlo}

Avec 2\,000 tirages $w'\sim\text{Dirichlet}(50\,w)$, la probabilité d'être classée première sous TOPSIS est
de 47,6\,\% pour M2, 42,5\,\% pour M4, 9,3\,\% pour M1 et 0,8\,\% pour M3 ; la corrélation de Spearman
moyenne avec le classement de référence est de 0,73. Ce résultat confirme l'analyse précédente : M2 et M4
sont pratiquement à égalité sous TOPSIS, et une petite incertitude sur les poids suffit à les départager.

\section{Recommandation}

Au vu de la concordance de cinq méthodes sur six, du consensus de Borda et de la fragilité du classement
TOPSIS, la recommandation la plus défendable est \textbf{M1}, en signalant au décideur que le résultat
dépend sensiblement du poids attribué au coût et au style. Cette démarche --- plusieurs méthodes,
comparaison des rangs, sensibilité, justification --- est celle que la plateforme rend systématique.

% =====================================================================
\chapter{Tests et assurance qualité}\label{ch:tests}
% =====================================================================

\section{Tests automatisés}

Le fichier \texttt{tests/test\_course\_examples.py} contient quinze tests, exécutables avec
\texttt{python -m pytest tests} ou directement avec \texttt{python tests/test\_course\_examples.py}
(aucune dépendance supplémentaire). Ils couvrent :
\begin{itemize}
  \item les valeurs publiées du cours : $\xi^*$ BWM, Spearman (0,77 ; 0,94 ; 0,60), tableau ROC des trois
        cas, évaluation floue d'une alternative ;
  \item les décisions des exercices : AHP (voiture 1, $CR<0{,}1$), WASPAS ($A_4$), TOPSIS (rangs
        $[3,1,4,2]$), VIKOR ($S$, $Q$, compromis $A_2$), CRITIC ($C_4$ dominant), Chen
        ($A_2\succ A_3\succ A_1$) ;
  \item des propriétés mathématiques : AHP exact sur une matrice parfaitement cohérente ($CR=0$), somme
        nulle des flux nets PROMETHEE, idempotence de IFWA, conservation des rapports de poids lors d'une
        perturbation, utilités MAUT.
\end{itemize}
Les quinze tests passent.

\section{Tests de l'interface}

L'ensemble des pages a été exécuté avec le module de test de Streamlit (\texttt{streamlit.testing.v1.AppTest}) :
navigation dans les sept pages, chargement de chacun des cinq exercices dans chaque page, exécution des six
méthodes de classement, des quatre modes d'évaluation absolue, des dix études de cas, changement du nombre
de décideurs dans les méthodes floues, modèles BWM linéaire et non linéaire, sensibilité et Monte Carlo
pour plusieurs méthodes. Aucune exception n'a été levée. Le serveur Streamlit démarre et répond
correctement.

\section{Contrôles de cohérence à l'exécution}

Le module \texttt{utils/validation.py} bloque ou signale : matrice non numérique ou incomplète, moins de
deux alternatives, noms dupliqués, critère constant, valeurs non positives incompatibles avec WPM ou
l'entropie, matrice AHP non réciproque ou hors échelle de Saaty, incohérences BWM ($a_{BB}\ne1$, maximum
de BO ne portant pas sur $W$, $a_{BW}$ différent dans les deux vecteurs), TFN non ordonnés, IFN avec
$\mu+\nu>1$, poids négatifs (et renormalisation automatique si $\sum w_j\ne1$).

% =====================================================================
\chapter{Limites et perspectives}\label{ch:limites}
% =====================================================================

\section{Limites}
\begin{itemize}
  \item \textbf{ELECTRE III} n'est pas implémentée ; seul PROMETHEE représente les méthodes de
        surclassement. Les diapositives correspondantes n'en détaillant pas la formulation, PROMETHEE suit
        la formulation standard de Brans.
  \item Les méthodes floues sont limitées aux nombres triangulaires et intuitionnistes ; les ensembles
        flous pythagoriciens, q-rung orthopair, picturaux et hésitants présentés dans le cours ne sont pas
        encore exploités dans une méthode de classement.
  \item Le MCDM de groupe n'est disponible que dans les méthodes floues ; les méthodes \emph{crisp}
        supposent un décideur unique.
  \item Les critères sont supposés indépendants (pas d'ANP).
  \item L'exemple de TOPSIS flou intuitionniste et l'exemple MAUT sont illustratifs, faute de données
        chiffrées dans le support.
\end{itemize}

\section{Perspectives}
\begin{itemize}
  \item Ajouter ELECTRE III (indices de concordance, discordance, crédibilité et distillation).
  \item Étendre TOPSIS aux ensembles flous pythagoriciens et q-ROF, et ajouter AHP et BWM flous.
  \item Intégrer les scénarios \og MCDM et IA \fg{} du cours : apprentissage d'un modèle sur les scores MCDM
        et mise à jour des poids par SHAP ($w_j^{new}=\alpha w_j^{AHP}+(1-\alpha)|\phi_j|/\sum|\phi_j|$),
        MCDM dynamique alimenté par des prévisions de séries temporelles, extraction de scores à partir de
        textes par un LLM.
  \item Générer automatiquement un rapport de décision (PDF) reprenant les étapes et la justification.
\end{itemize}

% =====================================================================
\chapter*{Conclusion générale}
\addcontentsline{toc}{chapter}{Conclusion générale}
% =====================================================================

Nous avons conçu et réalisé une plateforme d'aide à la décision multicritère qui couvre l'ensemble du
processus présenté dans le cours : formalisation du problème, pondération des critères par des méthodes
subjectives, objectives ou hybrides, classement par des méthodes compensatoires et de surclassement,
évaluation absolue d'une alternative isolée, traitement de l'imprécision par la logique floue et analyse de
robustesse.

Une architecture en couches sépare le moteur de calcul, entièrement testable, de l'interface Streamlit, qui
expose chaque étape intermédiaire avec les formules du cours. La confrontation systématique aux exercices
des diapositives montre que la plateforme reproduit les résultats publiés ; le seul écart relevé (poids BWM
de l'exemple du fournisseur) a été analysé et attribué au support lui-même, la valeur optimale de $\xi$ étant
retrouvée et l'optimum du modèle étant unique.

Au-delà du calcul, l'étude comparative du choix d'une voiture montre l'intérêt pratique de l'outil : cinq
méthodes sur six désignent la même alternative, TOPSIS s'en écarte et l'analyse de sensibilité révèle que
son choix est fragile. La plateforme permet ainsi d'appliquer concrètement les bonnes pratiques du cours :
ne jamais se fier à une seule méthode, mesurer l'accord entre classements, tester la stabilité face aux
poids et justifier explicitement la décision.

% =====================================================================
\appendix
% =====================================================================

\chapter{Installation et lancement}

\begin{lstlisting}[language=bash,numbers=none,caption={Installation et lancement}]
cd mcdm_app
pip install -r requirements.txt
streamlit run app.py              # ouvre http://localhost:8501
python tests/test_course_examples.py
\end{lstlisting}

Dépendances : \texttt{streamlit>=1.40}, \texttt{numpy>=1.26}, \texttt{scipy>=1.11},
\texttt{pandas>=2.1}, \texttt{plotly>=5.20}.

\chapter{Format d'import CSV}

La première colonne contient les noms des alternatives, la première ligne les noms des critères. Une ligne
optionnelle nommée \texttt{type} indique le sens de chaque critère (\texttt{+} ou \texttt{-}) ; à défaut,
tous les critères sont considérés comme positifs. Le séparateur (virgule ou point-virgule) est détecté
automatiquement. L'export produit le même format.

\begin{lstlisting}[language={},numbers=none,caption={Exemple de fichier CSV}]
,Style,Fiabilite,Economie carburant,Cout
M1,7,9,9,8
M2,8,7,8,7
M3,9,6,8,9
M4,6,7,8,6
type,+,+,+,-
\end{lstlisting}

\chapter{Correspondance cours -- code}

\begin{table}[H]
\centering
\caption{Correspondance entre les méthodes du cours et les fonctions}\label{tab:corresp}
\begin{tabular}{lll}
\toprule
\textbf{Méthode} & \textbf{Pages du cours} & \textbf{Fonction}\\
\midrule
AHP & 23--34 & \texttt{weighting.ahp}, \texttt{weighting.ahp\_rank}\\
BWM & 35--43 & \texttt{weighting.bwm}, \texttt{weighting.bwm\_nonlinear}\\
Entropie & 44--45 & \texttt{weighting.entropy}\\
CRITIC & 46--51 & \texttt{weighting.critic}\\
Hybride & 22, 92 & \texttt{weighting.hybrid}\\
WSM / WPM / WASPAS & 53--59 & \texttt{crisp\_ranking.wsm / wpm / waspas}\\
TOPSIS & 60--65 & \texttt{crisp\_ranking.topsis}\\
VIKOR & 66--72 & \texttt{crisp\_ranking.vikor}\\
PROMETHEE & 74 & \texttt{outranking.promethee}\\
Spearman & 77--79 & \texttt{sensitivity.spearman\_matrix}\\
MAUT, seuil & 81--84 & \texttt{single\_alternative.maut}, \texttt{decide}\\
ROC, Youden & 85--88 & \texttt{single\_alternative.roc\_table}, \texttt{best\_thresholds}\\
Distance à l'idéal & 89 & \texttt{single\_alternative.distance\_to\_ideal}\\
TFN, une alternative & 134 & \texttt{single\_alternative.fuzzy\_single}\\
TOPSIS flou (Chen) & 135--142 & \texttt{fuzzy\_mcdm.fuzzy\_topsis\_chen}\\
IFWA, TOPSIS intuitionniste & 116, 143--150 & \texttt{fuzzy\_mcdm.ifwa}, \texttt{if\_topsis}\\
\bottomrule
\end{tabular}
\end{table}

\end{document}
